\documentclass[11pt]{article}
\usepackage[spanish]{babel}
\usepackage[utf8]{inputenc}
\usepackage[T1]{fontenc}
\usepackage{geometry}
\usepackage{booktabs}
\usepackage{graphicx}
\usepackage{hyperref}
\geometry{margin=2.2cm}
\title{Reporte del sistema actual de Gemelo Digital SIMION}
\author{Versión interna para el grupo de trabajo}
\date{2026-07-01}

\begin{document}
\maketitle

\section{Resumen ejecutivo}
Construimos un gemelo digital del beamline de SIMION: un emulador que recibe ocho voltajes libres
$(V3,V6,V9,V10,V11,V12,V15,V18)$ y predice métricas del haz. El sistema actual usa un ensemble de
redes neuronales MLP diferenciables. Además de predicción directa, permite diseño inverso por gradiente:
dada una consigna de transmisión y ángulo, busca voltajes que deberían cumplirla y luego los valida en SIMION.

\section{Entradas y salidas}
Entradas: ocho voltajes libres del reto. Salidas principales: transmisión activa, fracción forward,
ángulo medio de impacto, dispersión angular y residual de $v_z$ sigma sobre una base física.

\section{Arquitectura}
El modelo es un ensemble de MLPs multitarea con capas ocultas 128--128--64, activación GELU,
dropout 0.03, AdamW y pérdida log-cosh ponderada. Las entradas y salidas se escalan antes del entrenamiento.

\section{Datos y desempeño}
Dataset base: 3272 filas. Filas con residual de velocidad observado: 1702.
MAE global de transmisión: 3.02 puntos porcentuales. MAE high-contact:
1.48 puntos porcentuales. MAE high-contact de theta mean:
0.117 grados. MAE high-contact de theta sigma: 0.116 grados.

\section{Control inverso}
El diseño inverso congela la red y optimiza los voltajes de entrada por backpropagation. Los barridos
validados muestran control de transmisión aproximadamente desde 20\% hasta 100\%, con degradación angular
en baja transmisión.

\section{Rúbrica del reto}
\begin{itemize}
\item C, dirección: soluciones de alta transmisión validadas.
\item A, exactitud: error bajo en región informativa high-contact.
\item I, consigna inversa: barridos inversos validados en SIMION.
\item E, honestidad: incertidumbre por ensemble y distinción de regiones confiables/no confiables.
\item D, diferenciabilidad: modelo PyTorch diferenciable.
\item F, eficiencia: loops presupuestados para estimar número mínimo de evaluaciones.
\item G, admisibilidad: salidas físicas recortadas y métricas calculadas sobre impactos válidos.
\end{itemize}

\section{Archivos clave}
\begin{itemize}
\item \texttt{dt/train\_dt.py}: entrenamiento.
\item \texttt{dt/predict\_dt.py}: predicción directa.
\item \texttt{dt/inverse\_design.py}: diseño inverso.
\item \texttt{dt/budget\_simion\_loop.py}: entrenamiento autónomo con presupuesto.
\item \texttt{dt/inverse\_active\_loop.py}: aprendizaje activo por fallos de consigna.
\item \texttt{dt/EXPERIMENT\_LOG.md}: bitácora técnica.
\end{itemize}

\end{document}
